\documentclass[a4paper,10pt]{report}\input{../0.Préambule/Préambule_cours_prof}\begin{document}

\pagestyle{empty}
\newpage
\noindent NOM, Prénom et classe : \dotfill

\section*{Les nombres entiers}
\addcontentsline{toc}{section}{Les nombres entiers}

\noindent \textit{L'usage de la calculatrice est \textbf{interdite}.} \hfill \textit{Durée : 55 minutes}

\medskip
\paragraph{Exercice 1}: \hfill \textbf{4 points}

\medskip
\noindent \begin{minipage}{12cm}
Dans la division ci-contre :

\medskip
$185$ est le \dotfill \medskip

$4$ est le  \dotfill \medskip

$39$ est le \dotfill \medskip

$2$ est le \dotfill \medskip
\end{minipage}
\begin{minipage}{5cm}
\begin{center}
\opdiv[displayintermediary=all,voperation=top,maxdivstep=2,dividendbridge]{158}{4}
\end{center}
\end{minipage}

\medskip
\paragraph{Exercice 2}: \hfill \textbf{4 points}

\medskip
\noindent Complète le tableau par oui ou non.

\medskip
\noindent \renewcommand{\arraystretch}{1.8}
\begin{tabularx}{\linewidth}{|>{\columncolor{lightgray!50}\centering\arraybackslash}X|*{3}{>{\centering\arraybackslash}X|}}
\hline
\rowcolor{lightgray!50} Le nombre est-il divisible par ... & $2$ & $5$ & $9$ \\\hline
$508$ & & & \\\hline
$990$ & & & \\\hline
$505$ & & & \\\hline
$\np{2 664}$ & & & \\\hline
$\np{1 505}$ & & & \\\hline
$\np{34 452}$ & & & \\\hline
\end{tabularx}

\medskip
\paragraph{Exercice 3}: \hfill \textbf{4 points}

\medskip
On donne les égalités : $$415=7 \times 59+2 \quad \text{ et } \quad 56 \times  57=3192$$

Sans effectuer de calculs donne le quotient et le reste des divisions euclidiennes suivantes. On fera tout particulièrement attention à la rédaction des réponses.
\begin{enumerate}
\renewcommand{\arraystretch}{1}
\begin{tabularx}{\linewidth}{*{4}{>{\arraybackslash}X}}
\item $415$ par $7$&
\item $415$ par $59$&
\item $3192$ par $56$&
\item $3192$ par $57$\\
\end{tabularx}
\end{enumerate} 

\noindent\grandscarreaux{\linewidth}{6.4cm}

\medskip
\paragraph{Exercice 4}: \hfill \textbf{6 points}

\medskip
\begin{enumerate}
\item Convertis chaque durée en minutes.
\begin{enumerate}
\renewcommand{\arraystretch}{1}
\begin{tabularx}{\linewidth}{*{2}{>{\arraybackslash}X}}
\item $15$h$07$min&
\item $16$h $17$min\\
\end{tabularx}
\end{enumerate}

\noindent\grandscarreaux{\linewidth}{3.2cm}
\item Convertis en heures et minutes.
\begin{enumerate}
\renewcommand{\arraystretch}{1}
\begin{tabularx}{\linewidth}{*{2}{>{\arraybackslash}X}}
\item $375$min&
\item $369$min\\
\end{tabularx}
\end{enumerate}

\noindent\grandscarreaux{\linewidth}{4cm}
\item Convertis en heures, minutes et secondes.
\begin{enumerate}
\renewcommand{\arraystretch}{1}
\begin{tabularx}{\linewidth}{*{2}{>{\arraybackslash}X}}
\item $\np{16 000}$s&
\item $\np{25 000}$s\\
\end{tabularx}
\end{enumerate}

\noindent\grandscarreaux{\linewidth}{4cm}
\end{enumerate}

\medskip
\paragraph{Exercice 5}: \hfill \textbf{3 points}

\medskip
Un fleuriste achète $100$ roses à $0,80$~\euro . Avec toutes ces roses, il fait autant de bouquets de $6$ roses qu'il peut. Il décide de vendre ses bouquets pour $15$~\euro . Le restant des roses est vendu à $1,5$\euro ~l'unité.

\medskip
\begin{enumerate}
\item Combien le fleuriste fera-t-il de bouquets ?
\item Combien va gagner le fleuriste si il vend tout son stock ?
\end{enumerate}

\noindent\grandscarreaux{\linewidth}{4.8cm}
\end{document}